%!TEX root = ../DigSig2.tex
\section{Digital Filter Realizations\buch{Chapter 7}}

\subsection{Direct Form\buchSeite{265}}
\label{sec:directform}
Also called \emph{direct form I}. Consider a second-order filter with
transfer function
\begin{equation*}
	H(z) = \frac{N(z)}{D(z)}
		 = \frac{b_0 + b_1z^{-1} + b_2z^{-2}}{1 + a_1z^{-1} + a_2z^{-2}}
\end{equation*}
with the I/O difference equation
\begin{equation*}
	y_n = -a_1 y_{n-1} - a_2 y_{n-2} + b_0 x_n + b_1 x_{n-1} + b_2 x_{n-2}
\end{equation*}\\

This results in the following block diagram: \\

\begin{center}
	\input{tikz/DirectForm}
\end{center}

and following sample processing algorithm:
\begin{lstlisting}[mathescape]
for $\text{each input sample}$ x do:
  v$_0$ = x
  w$_0$ = -a$_1$w$_1$-a$_2$w$_2$+b$_0$v$_0$+b$_1$v$_1$+b$_2$v$_2$
  y = w$_0$
  v$_2$ = v$_1$, w$_2$ = w$_1$
  v$_1$ = v$_0$, w$_1$ = w$_0$
\end{lstlisting}

\paragraph{Filter Order}

The order of a digital filter is equal to the highest power of $z^{-1}$
appearing in the denominator polynomial.

\[
H(z)=\frac{N(z)}{D(z)}
=\frac{b_0+b_1z^{-1}+\cdots+b_Lz^{-L}}
{1+a_1z^{-1}+\cdots+a_Mz^{-M}}
\]

The filter order is therefore

\[
\boxed{\text{Order}=M.}
\]

The numerator may have a lower order than the denominator.

\subsection{Canonical Form\buchSeite{271}}
\label{sec:canonicalform}
The most popular form. Also called \emph{direct form II}. It can be obtained from
the \nameref{sec:directform} by splitting it up, introducing an intermediate signal $w(n)$ and reversing it.

\begin{align*}
	H(z) &= \frac{1}{D(z)} \cdot N(z) \\
	w(n) &= x(n) - a_1 w(n-1) - a_2 w(n-2) \\
	y(n) &= b_0 w(n) + b_1 w(n-1) + b_2 w(n-2) \\
\end{align*}\\

\begin{center}
	\input{tikz/CanonicalForm}
\end{center}

\begin{lstlisting}[mathescape]
for $\text{each input sample}$ x do:
  w$_0$ = x - a$_1$w$_1$ - a$_2$w$_2$
  y = b$_0$w$_0$ + b$_1$w$_1$ + b$_2$w$_2$
  w$_2$ = w$_1$
  w$_1$ = w$_0$
\end{lstlisting}


\subsection{Cascade Form\buchSeite{277}}
\label{sec:cascadeform}
For filters with real impulse responses (all $a_i$ and $b_i$ are real), the
transfer function can be written by cascaded 2nd order sections (SOS).

\begin{align*}
	H(z) = \prod_{i=0}^{K-1}H_i(z) = \prod_{i=0}^{K-1}\frac
	{b_{i0}+b_{i1}z^{-1}+b_{i2}z^{-2}}
	{1+a_{i1}z^{-1}+a_{i2}z^{-2}} &&
\end{align*}

Usually the SOS are implemented in the canonical form. A first order term is
implemented by setting the $z^{-2}$ term to zero. The advantage of the cascade
form is the efficient implementation of 2nd order sections in DSPs.

\paragraph{Block Diagram}

The $K$ second-order sections are connected in series:

\adjustbox{max width=0.9\linewidth}{
    \begin{tikzpicture}

\matrix (m) [row sep=2mm,column sep=7mm]
{
\node[dspnodeopen,dsp/label=above] (m00) {$x(n)$}; &
\node[draw,minimum width=1.8cm,minimum height=1cm] (m01)
{
$\begin{array}{c}
H_1(z)\\
\text{SOS 1}
\end{array}$
}; &
\node[coordinate] (m02) {}; &
\node[draw,minimum width=1.8cm,minimum height=1cm] (m03)
{
$\begin{array}{c}
H_2(z)\\
\text{SOS 2}
\end{array}$
}; &
\node[coordinate] (m04) {}; &
\node[coordinate] (m05) {$\cdots$}; &
\node[coordinate] (m06) {}; &
\node[draw,minimum width=1.8cm,minimum height=1cm] (m07)
{
$\begin{array}{c}
H_K(z)\\
\text{SOS K}
\end{array}$
}; &
\node[dspnodeopen,dsp/label=above] (m08) {$y(n)$};
\\
};

\begin{scope}[start chain]

\chainin(m00);

\chainin(m01)[join=by dspconn];

\chainin(m03)[join=by dspconn];

\chainin(m05)[join=by dspconn];

\chainin(m07)[join=by dspconn];

\chainin(m08)[join=by dspconn];

\end{scope}

\draw (m01)--node[midway,above]{$v_1(n)$}(m03);

\draw (m03)--node[midway,above]{$v_{2}(n)$}(m05);

\draw (m05)--node[midway,above]{$v_{K-2}(n)$}(m07);

\end{tikzpicture}
}

Each second-order section is implemented in the canonical form (Direct Form II):

\begin{center}
    \begin{tikzpicture}

\matrix (m) [row sep=2.5mm,column sep=5mm]
{

\node[dspnodeopen,dsp/label=above] (m00) {$v_i(n)$}; &
\node[dspadder] (m01) {}; &
\node[coordinate] (m02) {}; &
\node[coordinate] (m03) {}; &
\node[dspmixer,dsp/label=below] (m04) {$b_{i0}$}; &
\node[dspadder] (m05) {}; &
\node[dspnodeopen,dsp/label=above] (m06) {$v_{i+1}(n)$};
\\

\node[coordinate](m10){};&
\node[coordinate](m11){};&
\node[coordinate](m12){};&
\node[dspsquare](m13){$\z^{-1}$};&
\node[coordinate](m14){};&
\node[coordinate](m15){};&
\node[coordinate](m16){};
\\

\node[coordinate](m20){};&
\node[coordinate](m21){};&
\node[dspmixer,dsp/label=below](m22){$-a_{i1}$};&
\node[coordinate](m23){};&
\node[dspmixer,dsp/label=below](m24){$b_{i1}$};&
\node[coordinate](m25){};&
\node[coordinate](m26){};
\\

\node[coordinate](m30){};&
\node[coordinate](m31){};&
\node[coordinate](m32){};&
\node[dspsquare](m33){$\z^{-1}$};&
\node[coordinate](m34){};&
\node[coordinate](m35){};&
\node[coordinate](m36){};
\\

\node[coordinate](m40){};&
\node[coordinate](m41){};&
\node[dspmixer,dsp/label=below](m42){$-a_{i2}$};&
\node[coordinate](m43){};&
\node[dspmixer,dsp/label=below](m44){$b_{i2}$};&
\node[coordinate](m45){};&
\node[coordinate](m46){};

\\
};

%%%%%%%%%%%%%%%%%%%%%%%%%%%%%%%%%%%%%%%%%%

\begin{scope}[start chain]

\chainin(m00);

\chainin(m01)[join=by dspconn];

\chainin(m04)[join=by dspconn];

\chainin(m05)[join=by dspconn];

\chainin(m06)[join=by dspconn];

\end{scope}

%%%%%%%%%%%%%%%%%%%%%%%%%%%%%%%%%%%%%%%%%%

\begin{scope}[start chain]

\chainin(m03);

\chainin(m13)[join=by dspconn];

\chainin(m23)[join=by dspline];

\chainin(m33)[join=by dspconn];

\chainin(m43)[join=by dspline];

\end{scope}

%%%%%%%%%%%%%%%%%%%%%%%%%%%%%%%%%%%%%%%%%%

\foreach \i in {2,4}
{
\begin{scope}[start chain]

\chainin(m\i3);

\chainin(m\i2)[join=by dspconn];

\chainin(m01)[join=by dspconn];

\end{scope}
}

%%%%%%%%%%%%%%%%%%%%%%%%%%%%%%%%%%%%%%%%%%

\foreach \i in {2,4}
{
\begin{scope}[start chain]

\chainin(m\i3);

\chainin(m\i4)[join=by dspconn];

\chainin(m05)[join=by dspconn];

\end{scope}
}

%%%%%%%%%%%%%%%%%%%%%%%%%%%%%%%%%%%%%%%%%%

\draw (m01)--node[midway,above]{$w_{i,0}$}(m04);

\draw (m23)--node[midway,above]{$w_{i,1}$}(m24);

\draw (m43)--node[midway,above]{$w_{i,2}$}(m44);

\end{tikzpicture}
\end{center}

For a first-order section, set

\[
a_{i2}=0,
\qquad
b_{i2}=0.
\]

\subsection{Cascade $\Leftrightarrow$ Canonical\buchSeite{284}}
To find the \nameref{sec:cascadeform} of a filter, $D(z)$ and $N(z)$ are
factored into 2nd order sections.

\begin{align*}
D(z) &= 1 + a_1 z^{-1} + a_2 z^{-2} + \dots + a_M z^{-M} &\\
&=(1-p_1z^{-1})(1-p_2z^{-1}) \dots (1-p_Mz^{-1})&
\end{align*}

For real-valued roots, any two factors can be combined:
\begin{equation*}
	(1 - p_1 z^{-1}) (1 - p_2 z^{-1}) = (1 - (p_1 + p_2) z^{-1} + p_1 p_2 z^{-2})
\end{equation*}

Conjugate complex roots are paired (if the coefficients are real):
\begin{equation*}
	(1 - p_1 z^{-1}) (1 - p_1^* z^{-1}) = 1 - 2 \text{Re}(p_1) z^{-1} + |p_1|^2 z^{-2}
\end{equation*}

\subsection{Impulse Response of Periodic Filters}

Many transfer functions in this chapter contain a denominator of the form

\[
1-z^{-D}.
\]

Using the geometric series,

\[
\boxed{
\frac{1}{1-z^{-D}}
=
\sum_{k=0}^{\infty}z^{-kD}
=
1+z^{-D}+z^{-2D}+\cdots
}
\]

the impulse response can be obtained without using partial fraction
expansion.

If

\[
H(z)=
\frac{B(z)}
{1-z^{-D}},
\]

then

\[
h(n)
=
b(n)
*
\sum_{k=0}^{\infty}\delta(n-kD),
\]

which means that the FIR impulse response $b(n)$ is repeated every $D$
samples.

Example:

\[
H(z)=
\frac{1+2z^{-1}}
{1-z^{-4}}
\]

results in

\[
h=
[1,2,0,0,\;
1,2,0,0,\;
1,2,0,0,\ldots].
\]